% ══════════════════════════════════════════════════════════════════════
% CHAPTER 1: INTRODUCTION
% ══════════════════════════════════════════════════════════════════════

\chapter{Introduction}
\label{ch:introduction}

\section{Motivation and Context}
\label{sec:motivation}

\revised{Seismology is a data-intensive domain within Earth sciences. Seismic monitoring networks continuously record events, generating datasets that describe subsurface processes such as tectonic stress accumulation and magmatic fluid migration. Interpreting hypocentral data in near-real time is important for volcanic and seismic risk assessment.}

\revised{Traditionally, earthquake catalogues have been analyzed through two-dimensional maps, cross-sections, and time-series plots. These representations do not fully capture the four-dimensional spatial and temporal distribution of seismic processes. An earthquake sequence migrating along a fault at depth may appear as a diffuse group in map view, but reveal a planar structure when examined in three dimensions and animated over time. This limitation of conventional visualisation tools motivated the development of \textbf{Seismic Viewer}.}

\revised{The project was developed during an internship at the Istituto Nazionale di Geofisica e Vulcanologia (INGV). Researchers studying the Campi Flegrei caldera required a tool to display spatial distributions, temporal evolution, and structural relationships within seismic sequences in an interactive desktop application.}

\section{Problem Statement}
\label{sec:problem}

Despite the availability of numerous Geographic Information System (GIS) platforms and seismological software packages, a significant technological gap persists. Existing tools generally fall into one of two categories:

\begin{enumerate}[label=(\alph*)]
    \item \textbf{General-purpose GIS platforms} (e.g., QGIS, ArcGIS): These offer geospatial analysis capabilities but are not designed for rendering dynamic 3D point clouds with seismic metadata in real time. Their architectures prioritize cartographic precision over interactivity, leading to delays when handling datasets with tens of thousands of hypocentres. \revised{From a software architecture perspective, traditional GIS tools often use synchronous data processing and standard text formats like JSON. When handling large datasets with tens of thousands of points, this causes the user interface to freeze due to high memory usage and parsing delays.}

    \item \textbf{Specialised seismological software} (e.g., Seisan, GMT, ObsPy): These provide robust analytical functionality but typically produce static outputs, such as static figures for print rather than interactive explorations. The learning curve is steep, requiring command-line proficiency, and the visualisation paradigm remains fundamentally two-dimensional.
\end{enumerate}

Neither category adequately addresses the need for a tool that unifies \textit{real-time 3D rendering}, \textit{temporal animation}, \textit{machine learning-based clustering}, and \textit{structural geological analysis} within a single, accessible desktop application. Seismic Viewer was designed to fill this gap.

\section{Objectives and Scope}
\label{sec:objectives}

The primary objectives of this project are:

\begin{enumerate}
    \item \textbf{Interactive 4D Visualisation}: Develop a system capable of rendering seismic hypocentres in three spatial dimensions plus time, with smooth animations and real-time filtering.

    \item \textbf{Automated Structural Analysis}: Integrate spatial clustering algorithms (DBSCAN, K-Means) and plane fitting (RANSAC + SVD) to automatically identify seismogenic structures. \revised{In particular, solve the cluster label switching issue across consecutive 4D animation frames to keep colors consistent over time.}

    \item \textbf{Multi-Source Data Integration}: Support ingestion from local CSV files, the INGV GOSSIP online API, seismic tomography datasets (Vp/Vs ratios), and GNSS ground deformation data.

    \item \textbf{Performance at Scale}: Ensure sub-second interactive response times for large datasets, through \revised{a custom Float32 binary streaming protocol, an asynchronous multi-channel communication design (native bridge, HTTP REST, WebSockets)}, GPU-accelerated rendering, and efficient caching strategies.

    \item \textbf{Accessible User Experience}: Provide a modern, intuitive interface that does not require programming expertise, enabling geophysicists to focus on interpretation rather than tool operation.
\end{enumerate}

This report covers the development process of the application, from the initial design to testing and evaluation. The system was evaluated using earthquake data from the Campi Flegrei caldera. \revised{In addition to practical usability, this thesis explains the software engineering solutions used to connect Python scientific calculations with interactive 3D web graphics.} The following chapters present the domain background, the system architecture, the data processing pipeline, the visualisation modes, the algorithmic foundations, the development process, and the conclusions drawn from this work.

\section{Internship Context}
\label{sec:internship}

This work was conducted as part of a university internship. The project was developed collaboratively with Giada Alfonsi, a PhD student in geology at Roma Tre University, under the supervision of Daniele Bailo, a software engineer at the INGV.

An iterative development approach was adopted, progressively expanding the software's capabilities over the duration of the internship. Rather than following a rigid specification, each development cycle was shaped by the continuous feedback and domain expertise provided by Giada Alfonsi, Professor Guido Giordano (volcanologist and professor at Roma Tre University), researchers at the INGV Osservatorio Vesuviano, and Daniele Bailo. The primary phases of the project were as follows:

\begin{itemize}
    \item \textbf{Phase 1}: Initial prototype, basic 3D visualisation, and local CSV ingestion.
    \item \textbf{Phase 2}: Major architectural refactoring with the introduction of K-Means clustering, a GUI overhaul, online data import from INGV, and structural plane visualisation.
    \item \textbf{Phase 3}: Performance optimisation (binary streaming, cylinder-based clustering) and initial project documentation.
    \item \textbf{Phase 4}: Implementation of sub-clustering and auto-tuning of DBSCAN parameters via K-distance analysis.
    \item \textbf{Phase 5}: Integration of tomography and GNSS uplift layers, followed by the final modular refactoring of the codebase.
\end{itemize}

This collaborative and iterative feedback process proved essential for aligning the software with real analytical workflows, as several key features (such as the cylinder grid partitioning and the sub-clustering interface) emerged directly from practical requirements and discussions during development.

\section{Document Structure}
\label{sec:structure}

The remainder of this thesis is organised as follows:

\begin{description}
    \item[\Cref{ch:background}] reviews the seismological and technological background, surveys related work, and justifies the key technology choices.
    \item[\Cref{ch:architecture}] presents the system architecture, detailing the backend-frontend separation, communication protocols, and threading model.
    \item[\Cref{ch:data}] describes the data processing pipeline, including CSV ingestion, online import, geospatial computations, and performance optimisations.
    \item[\Cref{ch:visualisation}] documents the four visualisation modes and the integration of auxiliary data layers (tomography, uplift, station positions).
    \item[\Cref{ch:algorithms}] formalises the clustering and structural analysis algorithms, with mathematical derivations and implementation details.
    \item[\Cref{ch:development}] reflects on the development process, discusses key technical challenges, and presents the solutions adopted.
    \item[\Cref{ch:conclusions}] summarises the contributions, acknowledges limitations, and outlines directions for future work.
\end{description}
